% from ega2-1-fr.tex
\oldpage[II]{36}
\begin{corollary}[2.5.15]
\label{II.2.5.15-fr}
Supposons $S$ engendrée par $S_1$. Pour tout $S$-module gradué $M$ et
tout $n\in\mathbf{Z}$, on a
\[
\label{II.2.5.15.1-fr}
  \widetilde{M(n)}=\widetilde{M}(n)
\tag{2.5.15.1}
\]
à un isomorphisme canonique près.
\end{corollary}

\begin{proof}
Cela résulte des définitions (\hyperref[II.2.5.10.2-fr]{2.5.10.2}) et
(\hyperref[II.2.5.10.1-fr]{2.5.10.1}), de la prop.
(\hyperref[II.2.5.13-fr]{2.5.13}) et de l'existence d'un isomorphisme
canonique $M(n)\xrightarrow{\sim}M\otimes_S S(n)$ de degré $0$ qui, à tout
$z\in(M(n))_h=M_{n+h}$, fait correspondre
$z\otimes1\in M_{n+h}\otimes(S(n))_{-n}\subset(M\otimes_S S(n))_h$.
\end{proof}

\begin{env}[2.5.16]
\label{II.2.5.16-fr}
Désignons par $S'$ l'anneau gradué tel que $S'_0=\mathbf{Z}$,
$S'_n=S_n$ pour $n>0$. Alors, si $f\in S_d$ ($d>0$), on a
$(S(n))_{(f)}=(S'(n))_{(f)}$ pour tout $n\in\mathbf{Z}$, car un élément de
$(S'(n))_{(f)}$ est de la forme $x/f^k$ avec $x\in S'_{n+kd}$ ($k>0$), et
on peut toujours prendre $k$ tel que $n+kd\neq0$. Comme
$X=\operatorname{Proj}(S)$ et $X'=\operatorname{Proj}(S')$ s'identifient
canoniquement (\hyperref[II.2.4.7-fr]{2.4.7, (ii)}), on voit que pour tout
$n\in\mathbf{Z}$, $\mathcal{O}_X(n)$ et $\mathcal{O}_{X'}(n)$ sont images
l'un de l'autre par l'identification précédente.

Notons d'autre part que pour tout $d>0$ et tout $n\in\mathbf{Z}$, on a
\[
  (S^{(d)}(n))_h=S_{(n+h)d}=(S(nd))_{hd}
\]
pour $f\in S_d$, on a donc $(S^{(d)}(n))_{(f)}=(S(nd))_{(f)}$. On sait que
les schémas $X=\operatorname{Proj}(S)$ et
$X^{(d)}=\operatorname{Proj}(S^{(d)})$ s'identifient canoniquement
(\hyperref[II.2.4.7-fr]{2.4.7, (i)})~; ce qui précède montre que si la
$S_0$-algèbre $S^{(d)}$ est engendrée par $S_d$, $\mathcal{O}_X(nd)$ et
$\mathcal{O}_{X^{(d)}}(n)$ sont images l'un de l'autre par cette
identification, pour tout $n\in\mathbf{Z}$.
\end{env}

\begin{proposition}[2.5.17]
\label{II.2.5.17-fr}
Soit $d$ un entier $>0$, et soit
$U=\bigcup_{f\in S_d}D_+(f)$. La restriction à $U$ de l'homomorphisme
canonique
$\mathcal{O}_X(nd)\otimes_{\mathcal{O}_X}\mathcal{O}_X(-nd)
\to\mathcal{O}_X$ est un isomorphisme pour tout entier $n$.
\end{proposition}

\begin{proof}
En vertu de (\hyperref[II.2.5.16-fr]{2.5.16}), on peut se borner au cas où
$d=1$, et la conclusion résulte alors de la démonstration de
(\hyperref[II.2.5.13-fr]{2.5.13}).
\end{proof}

\subsection{$S$-module gradué associé à un faisceau sur
$\operatorname{Proj}(S)$.}
\label{subsection:II.2.6-fr}

\emph{On suppose dans tout ce numéro que l'idéal $S_+$ de $S$ est engendré
par l'ensemble $S_1$ des éléments homogènes de degré $1$.}

\begin{env}[2.6.1]
\label{II.2.6.1-fr}
Le $\mathcal{O}_X$-Module $\mathcal{O}_X(1)$ est alors
\emph{inversible} (\hyperref[II.2.5.9-fr]{2.5.9})~; on pose alors, pour tout
$\mathcal{O}_X$-Module $\mathcal{F}$
(\hyperref[0.5.4.6-fr]{0, 5.4.6}),
\[
\label{II.2.6.1.1-fr}
  \Gamma_*(\mathcal{F})
  =\Gamma_*(\mathcal{O}_X(1),\mathcal{F})
  =\bigoplus_{n\in\mathbf{Z}}\Gamma(X,\mathcal{F}(n))
\tag{2.6.1.1}
\]
compte tenu de (\hyperref[II.2.5.14.2-fr]{2.5.14.2}). Rappelons
(\hyperref[0.5.4.6-fr]{0, 5.4.6}) que $\Gamma_*(\mathcal{O}_X)$ est muni
d'une structure d'\emph{anneau gradué}, et $\Gamma_*(\mathcal{F})$ d'une
structure de
\emph{module gradué} sur $\Gamma_*(\mathcal{O}_X)$.

Puisque $\mathcal{O}_X(n)$ est localement libre, $\Gamma_*(\mathcal{F})$
est un foncteur covariant additif
\emph{exact à gauche} en $\mathcal{F}$~; en particulier, si $\mathcal{J}$
est un faisceau d'idéaux de $\mathcal{O}_X$, $\Gamma_*(\mathcal{J})$
s'identifie canoniquement à un
\emph{idéal gradué} de $\Gamma_*(\mathcal{O}_X)$.
\end{env}

\begin{env}[2.6.2]
\label{II.2.6.2-fr}
Soit $M$ un $S$-module gradué~; pour tout $f\in S_d$ ($d>0$),
$x\mapsto x/1$ est un homomorphisme de groupes abéliens
$M_0\to M_{(f)}$, et comme $M_{(f)}$ s'identifie canoni-

